% Generated draft: requires mathematical and source review.
\documentclass[11pt]{article}
\usepackage[margin=25mm]{geometry}
\usepackage{fontspec}
\setmainfont{Latin Modern Roman}
\usepackage{amsmath,amssymb,mathtools,unicode-math}
\setmathfont{Latin Modern Math}
\usepackage{enumitem,needspace,xurl,hyperref}
\setlength{\parindent}{0pt}
\setlength{\parskip}{5pt}
\setlength{\emergencystretch}{3em}
\newcommand{\N}{\mathbb{N}}
\newcommand{\Z}{\mathbb{Z}}
\newcommand{\Q}{\mathbb{Q}}
\newcommand{\R}{\mathbb{R}}
\newcommand{\C}{\mathbb{C}}
\newcommand{\F}{\mathbb{F}}
\newcommand{\A}{\mathbb{A}}
\newcommand{\PP}{\mathbb{P}}
\newcommand{\E}{\mathbb{E}}
\newcommand{\eps}{\varepsilon}
\newcommand{\dd}{\,\mathrm{d}}
\begin{document}
\section*{UnsolvedMath ID: 30001253}
\textbf{Source problem code:} catalogue entry
\subsection{Definitions and mathematical statement}
Let \(F_U\) and \(F_P\) be subfields of a common field \(F_{\wp}\), with compatible inclusions of \(F\), such that
\[
F=F_U\cap F_P\qquad\text{inside }F_{\wp}.
\]
The supplied source requires these to be patching fields associated with a curve. It does not specify the curve, its base, the construction of the fields, or the meanings of the indices \(U,P,\wp\). The question below is therefore conditional on a choice of such patching data; the displayed intersection identity alone does not specify the intended geometric setting.

For a linear algebraic group \(G\) over \(F\) and a field extension \(K/F\), write \(G(K)\) for its group of \(K\)-rational points. Regard \(G(F_U)\) and \(G(F_P)\) as subgroups of \(G(F_{\wp})\) via the given field inclusions. Say that \(G\) has the factorization property for these data if the multiplication map of sets
\[
\mu_G:G(F_U)\times G(F_P)\longrightarrow G(F_{\wp}),\qquad (g_U,g_P)\longmapsto g_Ug_P,
\]
is surjective. Equivalently, this means
\[
\forall g\in G(F_{\wp})\ \exists g_U\in G(F_U)\ \exists g_P\in G(F_P):\quad g=g_Ug_P.
\]
No uniqueness of the factors is required.

\textbf{Question.} For which linear algebraic groups \(G/F\) does this factorization property hold for the curve patching fields in the source? The extracted formulation imposes no connectedness or rationality hypothesis on \(G\). It does not clarify whether the desired characterization is relative to fixed patching data or must hold uniformly over all admissible curve patching data; this quantifier remains unresolved.
\subsection{Short English statement}
Determine which linear algebraic groups over \(F\) allow every point over \(F_{\wp}\) to be written, in the stated order, as a point over \(F_U\) times a point over \(F_P\). The geometric construction of the patching fields and the intended scope of quantification over those fields require clarification.
\subsection{Sources}
\begin{itemize}
\item Supplied local UnsolvedMath registry record, permanent ID \(30001253\), source problem code \texttt{catalogue entry}, titled \emph{Matrix Factorization over Curve Patching Fields}; collection: Oberwolfach Reports Open Problems. The record gives proposed year 2009 and supplies no individual proposer attribution.
\item \emph{Quadratic Forms and Linear Algebraic Groups} (2009), Oberwolfach Reports, pp.\ 1375--1426, DOI \texttt{10.4171/owr/2009/25}. Supplied report link: \url{https://doi.org/10.4171/owr/2009/25}.
\item Harbater--Hartmann--Krashen, \emph{Applications of patching}. Supplied literature link: \url{https://arxiv.org/abs/0809.4481}.
\item Supplied literature assessment dated \texttt{2026-08-21}: the record labels the problem \texttt{partially\_solved} and reports that patching gives factorization for rational and retract-rational connected groups, while the extension to arbitrary other linear algebraic groups remains unresolved. This is supplied dated curation derived from the OWR material, not an independently verified literature assessment; its status and cited claims require independent verification.
\item The supplied record labels the statement \texttt{corrected\_verified} and reports that its mathematical content was checked against the official Oberwolfach report, with notation, context, and question form normalized without changing the source-level claim. This verification description is supplied metadata, not a claim of source consultation or verification for this draft.
\end{itemize}
\end{document}
